\paragraph{Appendix A.1: Necessary free-probability definitions.}
There are many excellent books on free probability theory. In particular, we refer to the book [26] for more details than the brief introduction given here.
\paragraph{$C^*$-algebras.}
Let $\mathcal A$ be a unital $C^*$-algebra. For our purposes, we can think of this as an algebra of bounded operators on a complex Hilbert space which is self-adjoint ($a\in\mathcal A$ implies $a^*\in\mathcal A$), closed in the operator norm $\|\cdot\|$, and contains the identity ($1\in\mathcal A$). A faithful trace $\varphi$ on $\mathcal A$ is a continuous linear functional $\varphi:\mathcal A\to\mathbb C$ that is unital ($\varphi(1)=1$), positive $\varphi(aa^*)\geq0$, and $\varphi(aa^*)=0$ iff $a=0$.
The pair $(\mathcal A,\varphi)$ where $\mathcal A$ is a unital $C^*$-algebra and $\varphi$ is a faithful trace on $\mathcal A$ is called a $C^*$-probability space. Elements of $\mathcal A$ are called non-commutative random variables. An example of a finite-dimensional $C^*$-probability space is the class $(M_n(\mathbb C),\operatorname{tr}_n)$, which is the class of $n\times n$ complex matrices with the normalized trace functional defined as $\operatorname{tr}_n(M)=\frac1n\sum_{i=1}^nM_{ii}$. General $C^*$-probability spaces allow us to extend these definitions to infinite-dimensional operators, which are needed to define a non-commutative analog of independence called free independence. Faithfulness of the trace $\varphi$ then ensures that $\|a\|=\lim_{m\to\infty}\varphi((aa^*)^m)^{1/2m}$ (see [26, Proposition 3.17]). In particular, this allows one to compute the norm $\|\cdot\|$ by using the trace method and taking higher powers of the trace functional $\varphi$, as we will see later.
An illustrative example of an infinite-dimensional $C^*$-algebra is the following: let $G$ be an infinite discrete group with the identity element 1 and let $\ell_2(G)$ be the associated complex Hilbert space. The canonical orthonormal basis for $\ell_2(G)$ is given by the standard basis vectors $\{e_g:g\in G\}$. The left-regular representation of $G$ maps elements of $G$ to bounded linear operators on $\ell_2(G)$ and is defined by its action on the standard basis vectors,
\[\rho(g)e_h=e_{gh},\quad\forall h\in G.\]
The operator norm closure of the linear span of $\{\rho(g):g\in G\}$ forms a unital algebra $\mathcal A$ (called the reduced $C^*$-algebra of $G$) and together with the faithful trace $\varphi$ defined as
\[\varphi\left(\sum_g\alpha_g\rho(g)\right)=\alpha_1\]
gives a $C^*$-probability space.
\paragraph{Free Independence.}
Let $(\mathcal A,\varphi)$ be a $C^*$-probability space and let $\{\mathcal A_i\}_{i=1}^n$ be unital $*$-subalgebras of $\mathcal A$. They are said to be free (or freely independent) if for all $k\in[n]$, for all indices $i_1,\ldots,i_k\in[n]$, and for all $a_1\in\mathcal A_{i_1},\ldots,a_k\in\mathcal A_{i_k}$ satisfying $\varphi(a_1)=\cdots=\varphi(a_k)=0$, the joint free moment,
\[\varphi(a_1\cdots a_k)=0\]
whenever $j_1\neq j_2,j_2\neq j_3,\ldots,j_{k-1}\neq j_k$, that is, the free moments vanish when all the neighboring elements in the sequence $a_1,\ldots,a_k$ come from subalgebras with distinct indices, for example, $\varphi(a_1a_2a_1^*a_2^*a_3a_2)=0$.
Non-commutative random variables $a_1,\ldots,a_n\in(\mathcal A,\varphi)$ are said to be free if the subalgebras $\{\mathcal A_i\}_{i=1}^n$ are free, where $\mathcal A_i$ is the unital $*$-subalgebra generated by $a_i$ (the linear span of all monomials $a_i^{\epsilon_1}a_i^{\epsilon_2}\cdots a_i^{\epsilon_r}$ where $\epsilon_1,\ldots,\epsilon_r\in\{1,*\}$ and $r\in\mathbb N\cup\{0\}$). Note that the corresponding unital $C^*$-subalgebras obtained by taking the norm closure of each $\mathcal A_i$ are also freely independent in this case (see [26, Exercise 5.23]).
We remark that the set of free non-commutative random variables is an empty set if the underlying $C^*$-probability space is finite (for instance $(M_n(\mathbb C),\operatorname{tr}_n)$), so to find non-trivial examples one needs to work with infinite-dimensional $C^*$-probability spaces.
\paragraph{Free Haar Unitaries and Free Groups.}
Let $(\mathcal A,\varphi)$ be a $C^*$-probability space. An element $u\in\mathcal A$ is a Haar unitary if it is a unitary, i.e. $uu^*=u^*u=1$, and if $\varphi(u^k)=0$ for all non-zero integers $k$. A family $S=\{u_1,\ldots,u_n\}\in\mathcal A$ in a $C^*$-probability space $(\mathcal A,\varphi)$ is called a free Haar unitary family if each $u\in S$ is a Haar unitary and if $u_1,\ldots,u_n$ are free. For notational convenience, let us define $S^*=\{u_1^*,\ldots,u_n^*\}$ to be the set of corresponding adjoints.
One can give a concrete construction of free Haar unitaries in terms of the free group. The free group $F_n$ with a generating set $S$ of size $n$ is an infinite discrete group constructed as follows: a word is defined to be product of elements of $S\cup S^*$ with $\bot$ denoting the empty word that contains no symbols. A word is called reduced if it does not contain a sub-word of the form $gg^*$ or $g^*g$ for $g\in S$. Given a word that is not reduced, the process of repeatedly removing such sub-words until it becomes reduced is called reduction. The free group $F_n$ consists of all reduced words that can be built from the symbols in $S\cup S^*$ with the group operation being a product of words followed by reduction. The identity is the empty word $\bot$.
Let $\rho$ denote the left-regular representation of the free group $F_n$. Then, writing $S=\{s_1,\cdots,s_n\}$, the family $u_i=\rho(s_i)$ for $i\in[n]$, is a free Haar unitary family in the $C^*$-probability space given by the reduced $C^*$-algebra of the group $F_n$. This connection to the free group also allows us to give a very precise condition when the trace $\varphi$ evaluated on a non-commutative monomial in the $u_i$'s vanishes. For a $d$-tuple $\mathbf i=(i_1,\ldots,i_d)\in[m]^d$, let $u_{\mathbf i}$ denote the non-commutative monomial $u_{i_1}\cdots u_{i_d}$ and write $u_{\mathbf i}^*=(u_{\mathbf i})^*=u_{i_d}^*\cdots u_{i_1}^*$. Let $\mathbf i^1,\ldots,\mathbf i^t,\mathbf j^1,\ldots,\mathbf j^t$ each be a $d$-tuple in $[m]^d$ and consider the degree-$2td$ non-commutative monomial $w=u_{\mathbf i^1}u_{\mathbf j^1}^*u_{\mathbf i^2}u_{\mathbf j^2}^*\cdots u_{\mathbf i^t}u_{\mathbf j^t}^*$. Note that a degree-$2td$ monomial $w$ corresponds to an ordered $2td$-tuple of variables. To illustrate, if $t=1,m=3$ and $\mathbf i^1=(1,2,3)$ and $\mathbf j^1=(2,2,1)$, then $w=u_1u_2u_3(u_2u_2u_1)^*=u_1u_2u_3u_1^*u_2^*u_2^*$ and corresponds to the ordered tuple $(u_1,u_2,u_3,u_1^*,u_2^*,u_2^*)$. We can also interpret $w$ as a word in the free group by applying the reduction rules. Then the next proposition follows from the definitions of free independence and Haar unitaries.
\paragraph{Proposition 13.}
$\varphi(w)=1$ iff $w$ reduces to identity in the free group $F_n$, and $\varphi(w)=0$ otherwise.
